% Part: first-order-logic
% Chapter: introduction
% Section: sentences

\documentclass[../../../include/open-logic-section]{subfiles}

\begin{document}

\olfileid{fol}{int}{snt}

\olsection{\usetoken{P}{sentence}}

এখন !!{structure}s ও !!{formula}s-এর জন্য পরিতৃপ্তির (“এর মধ্যে সত্য”
হওয়ার) একটি সংজ্ঞার অন্তত রূপরেখা আমাদের আছে। কিন্তু তার জন্য এই বাড়তি
উপকরণটি—!!a{variable} আরোপ—দরকার হয়েছে, অথচ আমরা চেয়েছিলাম
!!{sentence}s-এর সংজ্ঞা। আরোপগুলি কীভাবে বাদ দেব এবং !!{sentence}s-ই বা
কী?

বিধিবদ্ধ যুক্তিবিদ্যার প্রথম পরিচয়ে !!{variable}s ও পরিমাণসূচকের সম্পর্ক
নিয়ে আলোচনার কথা সম্ভবত তোমার মনে আছে। পরিমাণসূচকের ঠিক পরেই সব সময়
!!a{variable} থাকে; তারপর !!{sentence}-এর যে অংশে ওই পরিমাণসূচকটি প্রযোজ্য
(তার “পরিসর”), সেখানে !!{variable}-টি ওই পরিমাণসূচক দ্বারা “বদ্ধ” বলে
বোঝা হয়। !!{formula}s-এ প্রতিটি !!{variable}-এর জন্য সংশ্লিষ্ট পরিমাণসূচক
থাকার শর্ত ছিল না; সংশ্লিষ্ট পরিমাণসূচকবিহীন !!{variable}s হলো “মুক্ত” বা
“অবদ্ধ”। যেসব !!{formula}s-এ কোনো মুক্ত !!{variable}s নেই, তাদের সকলকেই
আমরা !!{sentence}s বলব।

আবারও, !!a{formula}~$!A$-তে !!a{variable}-এর কোনো সংঘটন কখন বদ্ধ, কোন
পরিমাণসূচক তাকে বাঁধে এবং কখন সেটি মুক্ত—এই স্বজ্ঞাত ধারণাটি বোঝা কঠিন
নয়। এটি পরীক্ষার কোনো পদ্ধতি, হয়তো বন্ধনী গোনার পদ্ধতি, তুমি শিখে
থাকতে পারো। আমাদের একটি নির্ভুল সংজ্ঞার উপর জোর দিতে হবে—আর যেহেতু
!!{formula}s-কে আরোহের সাহায্যে সংজ্ঞায়িত করেছি, তাই !!a{formula}~$!A$-তে
!!a{variable}~$x$-এর মুক্ত ও বদ্ধ সংঘটনগুলিকেও আরোহের সাহায্যে সংজ্ঞায়িত
করা যায়। যেমন, আমাদের সরলীকৃত ভাষায় সংজ্ঞাটি এমন হতে পারে:
\begin{enumerate}
  \item $!A$ পরমাণু হলে তার মধ্যে $x$-এর সব সংঘটন মুক্ত (অর্থাৎ,
  $\Atom{\Obj P}{x}$-এ $x$-এর সংঘটনটি মুক্ত)।
  \item $!A$ যদি $\lnot !B$ রূপের হয়, তবে $\lnot !B$-তে~$x$-এর কোনো
  সংঘটন মুক্ত যদি এবং কেবল যদি~$!B$-তে~$x$-এর অনুরূপ সংঘটনটি মুক্ত
  (অর্থাৎ,~$!B$-তে চলরাশিগুলির যে সংঘটনগুলি মুক্ত, ঠিক সেগুলির অনুরূপ
  সংঘটনই~$\lnot !B$-তে মুক্ত)।
  \item $!A$ যদি $(!B \land !C)$ রূপের হয়, তবে $(!B \land !C)$-তে~$x$-এর
  কোনো সংঘটন মুক্ত যদি এবং কেবল যদি~$!B$ বা~$!C$-তে~$x$-এর অনুরূপ
  সংঘটনটি মুক্ত।
  \item $!A$ যদি $\lexists[x][!B]$ রূপের হয়, তবে~$!A$-তে $x$-এর কোনো
  সংঘটনই মুক্ত নয়; আর সেটি যদি $\lexists[y][!B]$ রূপের হয়, যেখানে
  $y$,~$x$-এর থেকে আলাদা !!{variable}, তবে $\lexists[y][!B]$-তে~$x$-এর
  কোনো সংঘটন মুক্ত যদি এবং কেবল যদি~$!B$-তে~$x$-এর অনুরূপ সংঘটনটি মুক্ত।
\end{enumerate}

চলরাশির মুক্ত ও বদ্ধ সংঘটনের নির্ভুল সংজ্ঞা পাওয়ার পর আমরা সহজেই বলতে
পারি: !!a{sentence} হলো মুক্ত !!{variable}s-এর সংঘটনবিহীন যেকোনো
!!{formula}।

\end{document}
